\subsection{Sprint 0}

Na Sprint 0, o foco principal foi entender o projeto, alinhar o problema apresentado pelos stakeholders e definir as tecnologias que seriam utilizadas. Como era minha primeira experiência na AGES, essa sprint teve um papel muito importante de ambientação. Eu ainda estava entendendo como funcionavam as cerimônias, a divisão de tarefas e a dinâmica de trabalhar em um projeto com pessoas de diferentes níveis de experiência.

Depois da definição inicial da stack, dediquei boa parte do tempo a estudar \textit{React Native} \cite{reactnative} e TypeScript \cite{typescript}, pois seria meu primeiro contato prático com essas tecnologias. Também participei da criação dos mockups no Figma, ficando responsável por telas de administrador. No começo tive dificuldade com a ferramenta, principalmente por não conhecer bem seus recursos e a forma de organizar elementos visuais, mas com a prática e com a ajuda do time consegui concluir as telas previstas.

Essa sprint me mostrou que o trabalho em grupo exigia mais do que apenas escrever código. Aprendi a dividir tarefas, acompanhar o andamento pelo processo Scrum e pedir ajuda quando necessário. Como AGES I, foi um primeiro contato importante com a rotina de desenvolvimento de software em equipe.
